\sezione{2}{Approccio bayesiano e regola di classificazione}{13–15}

È il \textbf{cuore della lezione}: il teorema di Bayes della sezione 1 diventa un algoritmo di classificazione.

\subsection{Il problema}
\begin{intuizione}
Abbiamo oggetti da classificare (per esempio persone da etichettare come maschio o femmina). Di ogni oggetto misuriamo alcune caratteristiche (peso, altezza): è il \key{pattern} $\x$. Dobbiamo decidere a quale \key{classe} appartiene.
\end{intuizione}
\begin{tecnica}
\begin{itemize}
  \item $\mathbf{V}$: spazio dei pattern $d$-dimensionali; $\x=(x^1,\dots,x^d)$ è un vettore di $d$ feature (con peso e altezza, $d=2$).
  \item $W=\{w_1,\dots,w_s\}$: le $s$ classi, \textbf{disgiunte}. Ogni pattern appartiene a una sola classe: le classi formano una \emph{partizione}, che serve per scrivere $p(\x)$ come somma.
\end{itemize}
\end{tecnica}

\begin{attenzione}[Il «se» della slide 13]
«Se tutte le distribuzioni in gioco sono note, l'approccio bayesiano è la soluzione \textsc{ottima}.» In pratica le distribuzioni \emph{non} sono note e vanno stimate dai dati. Il resto della lezione (parametrico, MLE, multinormale) serve proprio a stimarle.
\end{attenzione}

\subsection{I quattro ingredienti della formula}
\begin{center}
\rowcolors{2}{chiaro}{white}
\begin{tabularx}{\linewidth}{>{\centering\arraybackslash}m{1.6cm} >{\raggedright\arraybackslash}m{3.1cm} Y}
\rowcolor{navy}\thd{Simbolo} & \thd{Nome} & \thd{Domanda a cui risponde}\\
$P(w_i)$ & probabilità \textbf{a priori} (\emph{prior}) & Prima di guardare il pattern, quanto è probabile la classe $w_i$?\\
$p(\x\,|\,w_i)$ & densità \textbf{condizionale} (\emph{likelihood}) & Se la classe fosse $w_i$, quanto sarebbe «tipico» osservare proprio $\x$?\\
$p(\x)$ & densità \textbf{assoluta} (\emph{evidence}) & Quanto è comune $\x$ in generale, senza conoscere la classe?\\
$P(w_i\,|\,\x)$ & probabilità \textbf{a posteriori} (\emph{posterior}) & Dopo aver visto $\x$, quanto è probabile la classe $w_i$?\\
\end{tabularx}
\end{center}

\begin{attenzione}[Due trappole]
\textbf{1.} $p(\x|w_i)\neq P(w_i|\x)$: sono domande inverse («tra chi ha la carie, quanti hanno mal di denti?» contro «tra chi ha mal di denti, quanti hanno la carie?»). Bayes serve proprio a \emph{invertire} il condizionamento. Dal training set è facile stimare $p(\x|w_i)$ (si guardano gli esempi di una classe), ma per classificare serve $P(w_i|\x)$.

\textbf{2.} $P$ maiuscola vs $p$ minuscola: le classi sono discrete e hanno vere \emph{probabilità}; i pattern sono continui e hanno \emph{densità}. Nella formula di Bayes le densità stanno sia sopra sia sotto la frazione: il fattore di scala si semplifica e $P(w_i|\x)$ è una vera probabilità.
\end{attenzione}

\subsection{La densità assoluta: legge della probabilità totale}
\begin{formula}
$\displaystyle p(\x)=\sum_{i=1}^{s} p(\x\,|\,w_i)\,P(w_i),\qquad\qquad \sum_{i=1}^{s}P(w_i)=1$
\end{formula}
\begin{intuizione}
Il pattern $\x$ può essere stato generato da una qualsiasi delle classi: si sommano i contributi di ogni classe, ciascuno pesato per quanto è comune la classe. «Quante persone alte 1,68\,m ci sono? Quelle tra i maschi $\times$ \% maschi + quelle tra le femmine $\times$ \% femmine.» $p(\x)$ è il \key{fattore di normalizzazione} che fa sommare a 1 le posteriori.
\end{intuizione}

\subsection{Il teorema di Bayes applicato}
\begin{formulat}{Teorema di Bayes}
$\displaystyle P(w_i\,|\,\x)=\frac{p(\x\,|\,w_i)\;P(w_i)}{p(\x)}
\qquad\qquad \text{posterior}=\frac{\text{likelihood}\times\text{prior}}{\text{evidence}}$
\end{formulat}
Derivazione: $p(\x,w_i)=P(w_i|\x)\,p(\x)=p(\x|w_i)\,P(w_i)$; si divide per $p(\x)$. Bayes è un \key{aggiornamento della credenza}: si parte dal prior, si osserva il dato, si moltiplica per quanto il dato è compatibile con la classe e si normalizza.

\subsection{La regola di classificazione (MAP)}
\begin{formulat}{Regola di Bayes (Maximum A Posteriori)}
$\displaystyle b=\argmax_{i=1,\dots,s} P(w_i\,|\,\x)\;=\;\argmax_{i=1,\dots,s}\; p(\x\,|\,w_i)\,P(w_i)$
\end{formulat}
\begin{tecnica}[Perché si può togliere $p(\x)$]
$p(\x)$ è uguale per tutte le classi: dividere tutti i candidati per lo stesso numero positivo non cambia chi vince. Per la \emph{decisione} conta solo il prodotto likelihood $\times$ prior; $p(\x)$ serve solo per il \emph{valore} della posteriori, cioè la confidenza (sezione 12).
\end{tecnica}
Due casi particolari:
\begin{itemize}
  \item \textbf{classi equiprobabili}: i prior si semplificano e decide la sola likelihood (regola \key{Maximum Likelihood});
  \item \textbf{likelihood uguali} (il pattern non dà informazione): decide il prior, cioè si sceglie sempre la classe più frequente.
\end{itemize}

\begin{esempio}[Il prior può ribaltare la decisione]
Due classi: $w_1$ = malattia rara, $w_2$ = sano. Si osserva un sintomo $\x$.
\begin{center}
\begin{tabular}{lccc}
\toprule
& prior $P(w_i)$ & likelihood $p(\x|w_i)$ & prodotto\\
\midrule
$w_1$ malato & $0.1$ & $\mathbf{0.5}$ & $0.05$\\
$w_2$ sano & $0.9$ & $0.1$ & $\mathbf{0.09}$\\
\bottomrule
\end{tabular}
\end{center}
$p(\x)=0.05+0.09=0.14$, quindi $P(w_1|\x)=0.05/0.14\approx\mathbf{0.36}$ e $P(w_2|\x)=0.09/0.14\approx\mathbf{0.64}$. La likelihood favorisce $w_1$ (il sintomo è 5 volte più tipico nei malati), ma i malati sono rari: vince \textbf{$w_2$}.
\end{esempio}
\fig[.8\linewidth]{prior_like.pdf}
\didascalia{Prior $\times$ likelihood $\Rightarrow$ posterior: la classe favorita dalla likelihood perde perché il suo prior è molto basso.}

\begin{mlbox}
La regola MAP non è arbitraria: è il classificatore che minimizza l'errore quando ogni sbaglio «costa» uguale (\emph{loss 0-1}). Se gli errori hanno costi diversi (in medicina un falso negativo è peggio di un falso positivo) si usa una versione con funzione di costo: si sceglie la classe con \emph{rischio atteso} minimo.
\end{mlbox}

\begin{appunti}{— Sezione 2}
\pag{13} Bayes è ottimo \textbf{solo se} le distribuzioni sono note; in pratica vanno stimate. Classi disgiunte = partizione. $p(\x|w_i)$ = likelihood (distribuzione dei pattern \emph{dentro} la classe); $P(w_i)$ = prior. $p(\x|w_i)\neq P(w_i|\x)$.
\pag{14} $p(\x)=\sum_i p(\x|w_i)P(w_i)$: probabilità totale, normalizzazione. Posterior = likelihood $\times$ prior / evidence. $P$ = probabilità (classi), $p$ = densità (pattern).
\pag{15} MAP: $b=\argmax_i P(w_i|\x)=\argmax_i p(\x|w_i)P(w_i)$. $p(\x)$ serve solo per la confidenza. Prior uguali → Maximum Likelihood. MAP = minimo errore con loss 0-1. Esempio: $0.05$ vs $0.09$, il prior ribalta la decisione.
\end{appunti}
